% Generated from locale/tr/content/normal-modal-logic/syntax-and-semantics/schemas.tex; do not hand-edit.


\olfileid[tr]{nml}{syn}{sch}

\olsection{Şemalar ve Geçerlilik}

\begin{defn}
  Bir \emph{şema}, bir !!{formula}s kümesidir; bu küme, $!C$ adlı bir kiplik
  !!{formula}{}ün bütün ve yalnızca bütün yerine koyma örneklerinden oluşur;
  yani
  \[
  \Setabs{!B}{\lexists[!D_1], \dots, \lexists[!D_n] \left(!B =
  \SSubst{!C}{\subst{!D_1}{p_1}, \dots, \subst{!D_n}{p_n}} \right) }.
  \]
  $!C$ adlı !!{formula}, şemanın \emph{karakteristik} !!{formula}{}dür;
  bu formül, !!{propositional variable}s{}nin yeniden adlandırılmasına kadar
  tektir. !!^a{formula}~$!A$, kümeye aitse şemanın bir \emph{örneği}dir.
\end{defn}

Bir şemayı, $!C$'nin atomik bileşenleri yerine `$!A$', `$!B$',~\dots,
koymakla elde edilen üstdilsel ifadeyle göstermek elverişlidir. Örneğin
`$!A$', `$!A \lif \Box!A$', `$!A \lif (!B \lif !A)$' ifadeleri şemaları
gösterir. Bunlar sırasıyla $p$, $p \lif \Box p$, $p \lif (q \lif p)$
karakteristik !!{formula}s{}ine karşılık gelir. `$!A$' şeması \emph{bütün}
!!{formula}s kümesini gösterir.

\begin{defn}
  Bir şema, bütün örnekleri bir modelde doğruysa ve ancak o zaman o modelde
  \emph{doğru}dur; her modelde doğruysa ve ancak o zaman \emph{geçerli}dir.
\end{defn}

\begin{prop}\ollabel{prop:Kvalid}
  Aşağıdaki \Ax{K} şeması geçerlidir:
  \begin{equation*}
  \Box(!A \lif !B) \lif (\Box !A \lif \Box !B). \tag{\Ax{K}}
  \end{equation*}
\end{prop}

\begin{proof}
  Şemanın bütün örneklerinin her modelin her dünyasında doğru olduğunu
  göstermeliyiz. Öyleyse $\mModel{M} = \tuple{W,R,V}$ ve $w \in W$ keyfî
  olsun. Bir koşullu önermenin bir dünyada doğru olduğunu göstermek için ön
  bileşenin doğru olduğunu varsayıp art bileşenin de doğru olduğunu
  gösteririz. Bu durumda $\mSat{M}{\Box(!A \lif !B)}[w]$ ve
  $\mSat{M}{\Box !A}[w]$ olsun. $\mSat{M}{\Box !B}[w]$ olduğunu
  göstermeliyiz. $Rww'$ olacak biçimde keyfî bir $w'$ alalım. İlk varsayım
  gereği $\mSat{M}{!A \lif !B}[w']$, ikinci varsayım gereği de
  $\mSat{M}{!A}[w']$ olur. Buradan $\mSat{M}{!B}[w']$ çıkar. $w'$ keyfî
  olduğundan $\mSat{M}{\Box !B}[w]$ olur.
\end{proof}

\begin{prop}\ollabel{prop:Dual-valid}
  Aşağıdaki \Dual{} şeması geçerlidir:
  \begin{equation*}
  \Diamond !A \liff \lnot\Box\lnot !A. \tag{\Dual}
  \end{equation*}
\end{prop}

\begin{proof}
  Alıştırma.
\end{proof}

\begin{prob}
  \olref[nml][syn][sch]{prop:Dual-valid} önermesini ispatlayın.
\end{prob}

\begin{prop}\ollabel{prop:soundMP}
  $!A$ ile $!A \lif !B$ bir modelin bir dünyasında doğruysa $!B$ de
  doğrudur. Dolayısıyla geçerli !!{formula}s modus ponens altında kapalıdır.
\end{prop}

\begin{prop}\ollabel{prop:valid-instances}
  !!^a{formula}~$!A$, ancak ve ancak bütün yerine koyma örnekleri geçerliyse
  geçerlidir. Başka bir deyişle bir şema, ancak ve ancak karakteristik
  !!{formula} geçerliyse geçerlidir.
\end{prop}

\begin{proof}
  ``Eğer'' yönü açıktır; çünkü $!A$ kendisinin bir yerine koyma örneğidir.

  ``Ancak'' yönünü ispatlamak için şunu gösterelim: $\mModel{M} =
  \tuple{W, R, V}$ bir kiplik modeli ve $!B \ident
  \SSubst{!A}{\subst{!D_1}{p_1},\dots,\subst{!D_n}{p_n}}$, $!A$'nın bir
  yerine koyma örneği olsun. $V'(p_i) =
  \Setabs{w}{\mSat{M}{!D_i}[w]}$ koyarak $\mModel{M'} = \tuple{W, R,
    V'}$ modelini tanımlayalım. O zaman her $w \in W$ için
  $\mSat{M}{!B}[w]$ ancak ve ancak $\mSat{M'}{!A}[w]$ olur. (İspatı
  alıştırma olarak bırakıyoruz.) Şimdi $!A$'nın geçerli, fakat $!A$'nın bir
  yerine koyma örneği olan $!B$'nin geçersiz olduğunu varsayalım. O zaman
  bir $\mModel{M} = \tuple{W, R, V}$ modeli ve bir $w \in W$ için
  $\mSat/{M}{!B}[w]$ olur. Fakat iddia gereği
  $\mSat/{M'}{!A}[w]$ olur; dolayısıyla $!A$ geçerli değildir. Bu bir
  çelişkidir.
\end{proof}

\begin{prob}
  \olref[nml][syn][sch]{prop:valid-instances} önermesinin ispatındaki
  ``ancak'' kısmında ileri sürülen iddiayı ispatlayın. (İpucu: $!A$ üzerinde
  tümevarım kullanın.)
\end{prob}

Bununla birlikte bir şemanın bir modelde doğru olmasının, ancak ve ancak
karakteristik formülünün o modelde doğru olması demek olmadığına dikkat
edin. Elbette ``ancak'' yönü geçerlidir: $!A$'nın her örneği
$\mModel{M}$ modelinde doğruysa $!A$'nın kendisi de $\mModel{M}$ modelinde
doğrudur. Fakat $!A$, $\mModel{M}$ modelinde doğruyken $!A$'nın bir örneği
$\mModel{M}$ modelinin bir dünyasında yanlış olabilir. Çok basit bir karşı
örnek için yalnızca bir $w$ dünyası bulunan ve $V(p) = \{w\}$ olan bir
modelde $p$'yi ele alın; böylece $p$, $w$'de doğrudur. Ancak $\lfalse$,
$p$'nin bir örneğidir ve $w$'de doğru değildir.

\begin{prob}
  Aşağıdaki !!{formula}s{}in hiçbirinin geçerli olmadığını gösterin:
  \begin{enumerate}
    \item[\Ax{D}:] \quad $\Box p \lif \Diamond p$;
    \item[\Ax{T}:] \quad $\Box p \lif p$;
    \item[\Ax{B}:] \quad $p \lif \Box\Diamond p$;
    \item[\Ax{4}:] \quad $\Box p \lif \Box \Box p$;
    \item[\Ax{5}:] \quad $\Diamond p \lif \Box \Diamond
      p$.
  \end{enumerate}
\end{prob}

\begin{table}[t]
    \centering
    \begin{tabular}{| l || l |}
      \hline
      {\emph{Geçerli Şemalar}} & {\emph{Geçersiz Şemalar}} \\
      \hline\hline
      $\Box(!A \lif !B) \lif (\Diamond !A \lif \Diamond !B)$
      & $\Box (!A \lor !B) \lif (\Box !A \lor \Box !B)$ \\
      $\Diamond (!A \lif !B) \lif (\Box !A \lif \Diamond
      !B)$
      & $(\Diamond !A \land \Diamond !B) \lif \Diamond (!A
      \land !B)$\\
      $\Box (!A \land !B) \liff (\Box !A \land \Box !B)$
      & $!A \lif \Box !A$ \\
      $\Box !A \lif \Box (!B \lif !A)$
      & $\Box \Diamond !A \lif !B$ \\
      $\lnot \Diamond !A \lif \Box (!A \lif !B)$
      & $\Box \Box !A \lif \Box !A$ \\
      $\Diamond (!A \lor !B) \liff (\Diamond !A \lor
      \Diamond !B)$
      & $\Box \Diamond !A \lif \Diamond \Box !A$. \\
      \hline
    \end{tabular}
    \caption{Geçerli ve geçersiz şemalar.}
    \ollabel{tab:valid-invalidSchemas}
\end{table}

\begin{prob}%\ollabel{ex:in/validSchemas}
  \olref[nml][syn][sch]{tab:valid-invalidSchemas} tablosunun ilk
  sütunundaki şemaların geçerli, ikinci sütunundakilerin ise geçersiz
  olduğunu ispatlayın.
\end{prob}

\begin{prob}
  Aşağıdaki şemaların geçerli mi, geçersiz mi olduğuna karar verin:
  \begin{enumerate}
  \item $(\Diamond !A \lif \Box !B) \lif (\Box !A \lif \Box
    !B)$;
  \item $\Diamond(!A \lif !B) \lor \Box(!B \lif !A)$.
  \end{enumerate}
\end{prob}

\begin{prob}
  Aşağıdaki şemaların her biri için, !!{formula}{}ün her örneğinin
  $\mModel{M}$ modelinde doğru olduğu bir $\mModel{M}$ modeli bulun:
  \begin{enumerate}
  \item $p \lif \Diamond\Diamond p$;
  \item $\Diamond p \lif \Box p$.
  \end{enumerate}
\end{prob}

